%% kapa: exam=1 round=36 year=2025 period=2 form=A subject=appraisal_law subject_ko=감정평가관계법규 seq=15 number=15
%% source: Q-Net 감정평가사 기출문제 (cst00309) artlSeq=5231525 file=raw/1008/5231525/2025년 제36회 감정평가사 1차시험 2교시 문제지 원본.pdf
%% conversion: kordoc 4.24.51 -> tools/gen1.py (auto); review=proofread-sonnet; flags=none
\begin{question}{15}
부동산 가격공시에 관한 법령상 개별공시지가에 관한 설명으로 옳지 않은 것은?
\begin{choices}
  \item 군수가 관계 행정기관의 장과 협의하여 개별공시지가를 결정ㆍ공시하기로 한 토지에 대해서는 개별공시지가를 결정ㆍ공시하여야 한다.
  \item 토지의 형질변경으로 ｢공간정보의 구축 및 관리 등에 관한 법률｣에 따른 지목변경이 된 토지는 개별공시지가 공시기준일을 다르게 할 수 있다.
  \item 시장은 개별공시지가의 결정ㆍ공시를 위하여 개별토지의 가격을 산정하는 경우 용도지역이 변경되는 토지를 검증 생략 대상 토지로 선정할 수 있다.
  \item 개별공시지가에 이의가 있는 자는 그 결정ㆍ공시일부터 30일 이내에 서면으로 시장ㆍ군수 또는 구청장에게 이의를 신청할 수 있다.
  \item 군수는 토지가격비준표의 적용에 오류가 있음을 발견한 때에는 지체 없이 이를 정정하여야 한다.
\end{choices}
\end{question}
\answer{3}
